\documentclass[11pt,a4paper]{article}

\usepackage[margin=2.6cm]{geometry}
\usepackage{amsmath,amssymb,amsthm}
\usepackage{bm}
\usepackage[colorlinks=true,linkcolor=blue,citecolor=blue,urlcolor=blue]{hyperref}

\newtheorem{remark}{Remark}

% --- Notation macros ---
\newcommand{\R}{\mathbb{R}}
\newcommand{\tr}{\operatorname{tr}}
\newcommand{\logdet}{\operatorname{log\,det}}
\newcommand{\dd}{\mathrm{d}}
\newcommand{\T}{\mathsf{T}}
\newcommand{\Gnoise}{\Gamma_{\mathrm{noise}}}
\newcommand{\Gprior}{\Gamma_{\mathrm{pr}}}
\newcommand{\Gpost}{\Gamma_{\mathrm{post}}}
\newcommand{\sensor}[1]{\bm{s}_{#1}}
\newcommand{\code}[1]{\texttt{#1}}

\title{Bayesian Optimization of Sensor Positions in\\
Magnetic Detection Electrical Impedance Tomography}
\author{Theory companion to \code{studies/optimal\_sensors\_bayesian\_approach/main.m}}
\date{\today}

\begin{document}
\maketitle

\begin{abstract}
This document describes the theory implemented in \code{main.m} for optimizing the positions of
magnetometers in magnetic detection electrical impedance tomography (MDEIT). The approach follows
the Bayesian optimal experiment design framework of Hyv\"onen, Sepp\"anen and Staboulis
[\emph{Optimizing electrode positions in electrical impedance tomography}, SIAM J.\ Appl.\ Math.,
arXiv:1404.7300], who optimized \emph{electrode} positions in EIT. Here the design variables are
instead the positions of external magnetic field sensors. The complete electrode model is
linearized around the prior mean conductivity, which yields an explicit Gaussian posterior
covariance. The sensor positions are then chosen to minimize either the trace
(A-optimality) or the log-determinant (D-optimality) of the posterior covariance, using analytic
gradients obtained from adjoint solutions and explicit derivatives of the Biot--Savart kernel.
\end{abstract}

\section{Problem setting}

Let $\Omega \subset \R^3$ be a conducting body with electrical conductivity $\sigma$, and let
$M \geq 1$ magnetometers be located at positions $\sensor{m} \in \R^3 \setminus \overline{\Omega}$,
$m = 1,\dots,M$. Currents are injected through $L$ fixed surface electrodes using $N$ stimulation
patterns. Each injected pattern drives an internal current density
$\bm{J} = -\sigma \nabla u$, whose magnetic flux density is measured by the sensors.

The inverse problem of MDEIT is to recover the (discretized) conductivity
$\sigma \in \R^{n_{\mathrm{el}}}$ from these magnetic measurements. The \emph{design} problem
considered here is: given a prior distribution for $\sigma$ and a noise model, place the sensors
$\{\sensor{m}\}$ so that the posterior distribution of $\sigma$ given the data is as
localized (informative) as possible.

Two features distinguish this setting from the EIT electrode-position problem of Hyv\"onen et al.:
\begin{enumerate}
  \item The electrodes are \emph{fixed}; only the contact-free magnetometers move. Consequently the
        forward PDE and its finite element (FE) system matrix do not depend on the design
        variables.
  \item The sensor positions enter the measurement map only through the smooth Biot--Savart
        kernel. The derivatives with respect to the design variables are therefore ordinary
        (explicit) derivatives of a volume integral; no shape derivatives of the boundary value
        problem are required. This is analytically and computationally simpler than the
        electrode-position case, which requires second-order Fr\'echet (shape) derivatives of the
        complete electrode model.
\end{enumerate}

\section{Forward model}

\subsection{Complete electrode model}

The electric potential $u$ inside $\Omega$ and the electrode potentials
$U = [U_1,\dots,U_L]^\T$ satisfy the complete electrode model (CEM):
\begin{equation}
\label{eq:cem}
\begin{aligned}
  \nabla \cdot (\sigma \nabla u) &= 0 && \text{in } \Omega, \\
  \sigma \, \partial_\nu u &= 0 && \text{on } \partial\Omega \setminus \cup_e E_e, \\
  u + z_e \, \sigma \, \partial_\nu u &= U_e && \text{on } E_e, \quad e = 1,\dots,L, \\
  \int_{E_e} \sigma \, \partial_\nu u \, \dd S &= I_e, && e = 1,\dots,L,
\end{aligned}
\end{equation}
where $z_e$ are the contact impedances and $I = [I_1,\dots,I_L]^\T$ the injected current pattern
(with $\sum_e I_e = 0$). Discretizing with piecewise linear finite elements yields the symmetric
block system
\begin{equation}
  \begin{pmatrix} A_c & A_e \\ A_e^\T & A_d \end{pmatrix}
  \begin{pmatrix} u \\ U \end{pmatrix}
  =
  \begin{pmatrix} 0 \\ I \end{pmatrix},
\end{equation}
where $A_c \in \R^{n_n \times n_n}$ contains the $\sigma$-weighted stiffness matrix plus contact
terms, and $A_e$, $A_d$ contain only contact-impedance terms (independent of $\sigma$).
Eliminating $U$ gives the Schur complement
\begin{equation}
\label{eq:schur}
  S(\sigma) \;=\; A_c - A_e A_d^{-1} A_e^\T,
\end{equation}
which is the symmetric positive (semi)definite operator used for all forward and adjoint solves
(function \code{M} in the code). Only $A_c$ depends on $\sigma$, so
\begin{equation}
\label{eq:dSdsigma}
  v^\T \frac{\partial S}{\partial \sigma_k}\, w
  \;=\; \int_{T_k} \nabla v \cdot \nabla w \, \dd V
  \;=\; |T_k| \, (\nabla v)_k \cdot (\nabla w)_k ,
\end{equation}
for any FE functions $v,w$, where $T_k$ is the $k$-th tetrahedron and the gradients are constant
per element for $P_1$ elements.

\subsection{Biot--Savart law and its discretization}

The magnetic flux density at sensor position $\sensor{m}$ generated by the internal current
density $\bm{J} = -\sigma\nabla u$ is
\begin{equation}
\label{eq:biotsavart}
  \bm{B}(\sensor{m})
  \;=\; \frac{\mu_0}{4\pi} \int_\Omega
        \bm{J}(\bm{r}) \times \frac{\sensor{m} - \bm{r}}{|\sensor{m} - \bm{r}|^3} \, \dd V .
\end{equation}
With $P_1$ elements, $\nabla u$ is constant on each element and is obtained from the sparse
gradient matrices $G_x, G_y, G_z \in \R^{n_{\mathrm{el}} \times n_n}$, i.e.\
$(\partial_x u)_k = (G_x u)_k$, etc. The geometric part of the kernel is collected in the dense
matrices $R_x, R_y, R_z \in \R^{M \times n_{\mathrm{el}}}$,
\begin{equation}
\label{eq:Rdef}
  (R_j)_{mk} \;=\; \int_{T_k} \frac{(\sensor{m} - \bm{r})_j}{|\sensor{m} - \bm{r}|^3}\, \dd V,
  \qquad j \in \{x,y,z\},
\end{equation}
evaluated with a 35-point quadrature rule on each tetrahedron
(\code{compute\_r\_matrices\_local}, \code{dRmkj\_dp}). Writing
$\Sigma = \operatorname{diag}(\sigma)$, expansion of the cross product in
\eqref{eq:biotsavart} yields the Cartesian components of the field at all sensors:
\begin{equation}
\label{eq:Cmatrices}
\begin{aligned}
  C_x(\sigma) &= -R_z \Sigma G_y + R_y \Sigma G_z, \\
  C_y(\sigma) &= -R_x \Sigma G_z + R_z \Sigma G_x, \\
  C_z(\sigma) &= -R_y \Sigma G_x + R_x \Sigma G_y,
\end{aligned}
\qquad
  B_j(\sensor{m}) = \frac{\mu_0}{4\pi}\, \big(C_j\, u\big)_m .
\end{equation}
Each sensor $m$ measures the projection of $\bm{B}$ onto its axes
$\hat{\bm{g}}^{(m)}_d$, $d = 1,2,3$, stored as $g_{m,d,j}$. The measurement operator for
axis $d$ is
\begin{equation}
\label{eq:Gamma}
  \Gamma_d(\sigma;\, \bm{s})
  \;=\; \frac{\mu_0}{4\pi}\,
        \big(\, g_{:,d,x} \odot C_x + g_{:,d,y} \odot C_y + g_{:,d,z} \odot C_z \,\big)
  \;\in\; \R^{M \times n_n},
\end{equation}
where $\odot$ denotes row-wise scaling (\code{compute\_gamma\_matrices\_local}). The measured
datum for sensor $m$, stimulation $l$ and axis $d$ is then simply
\begin{equation}
  B_{d,m}^{(l)} \;=\; \big(\Gamma_d\, u^{(l)}\big)_m,
  \qquad S(\sigma)\, u^{(l)} \;=\; b^{(l)} .
\end{equation}
Note that $\Gamma_d$ is \emph{linear} in $\sigma$ (through $\Sigma$) and depends on the sensor
positions only through the matrices $R_j$ in \eqref{eq:Rdef}.

For \emph{1-axis} MDEIT one fixed axis $d$ is measured, giving $MN$ data points; for
\emph{3-axis} MDEIT the data of all three axes are stacked, giving $3MN$ data points.

\section{Bayesian inversion and the linearized posterior}

Following the Bayesian paradigm, the discretized conductivity is a random vector with a Gaussian
prior, and the data are corrupted by additive Gaussian noise,
\begin{equation}
  \sigma \sim \mathcal{N}(\sigma_*,\, \Gprior),
  \qquad
  V \;=\; \mathcal{B}(\sigma;\, \bm{s}) + \varepsilon,
  \qquad
  \varepsilon \sim \mathcal{N}(0,\, \Gnoise),
\end{equation}
where $\mathcal{B}$ stacks all measurements and $\bm{s} = (\sensor{1},\dots,\sensor{M})$ collects
the design variables. To obtain a computationally tractable design criterion, the measurement map
is linearized around the prior mean,
\begin{equation}
  \mathcal{B}(\sigma) \;\approx\; \mathcal{B}(\sigma_*) + J\,(\sigma - \sigma_*),
  \qquad
  J = J(\sigma_*;\, \bm{s}) = \frac{\partial \mathcal{B}}{\partial \sigma}\Big|_{\sigma_*} .
\end{equation}
Under this linearization the posterior is Gaussian, with covariance
\begin{equation}
\label{eq:posterior}
  \Gpost \;=\; H^{-1},
  \qquad
  H \;=\; J^\T \Gnoise^{-1} J + \Gprior^{-1} .
\end{equation}
The essential observation for experiment design is that $\Gpost$ \emph{does not depend on the
measured data} --- only on the Jacobian $J$, hence on the sensor positions $\bm{s}$. The design
can therefore be optimized offline, before any measurement is taken.

\section{Optimality criteria}

\paragraph{A-optimality.}
Minimizing the expected mean-square error of the conditional-mean estimate (with unit weighting)
is equivalent to minimizing the trace of the posterior covariance:
\begin{equation}
\label{eq:aopt}
  \Phi_{\mathrm{A}}(\bm{s}) \;=\; \tr\!\big(\Gpost\big) \;=\; \tr\!\big(H^{-1}\big)
  \qquad \text{(\code{compute\_cost\_function\_a\_opt})}.
\end{equation}

\paragraph{D-optimality.}
Maximizing the expected information gain (the Kullback--Leibler divergence from the prior to the
posterior, marginalized over the data) is, for Gaussian linearized models, equivalent to
minimizing the log-determinant of the posterior covariance:
\begin{equation}
\label{eq:dopt}
  \Phi_{\mathrm{D}}(\bm{s}) \;=\; \logdet \Gpost \;=\; -\logdet H
  \qquad \text{(\code{compute\_cost\_function\_d\_opt})}.
\end{equation}
To avoid overflow/underflow, the determinant is evaluated through the Cholesky factorization
$H = LL^\T$:
\begin{equation}
  \logdet H \;=\; 2 \sum_i \log L_{ii} .
\end{equation}

\section{Gradients of the design criteria}

Let $p$ denote one Cartesian coordinate of one sensor position. Differentiating
$H = J^\T \Gnoise^{-1} J + \Gprior^{-1}$,
\begin{equation}
  \frac{\partial H}{\partial p}
  \;=\;
  \frac{\partial J}{\partial p}^{\!\T} \Gnoise^{-1} J
  \;+\;
  J^\T \Gnoise^{-1} \frac{\partial J}{\partial p} .
\end{equation}
Using $\dd\, \tr(H^{-1}) = -\tr(H^{-1}\, \dd H\, H^{-1})$ and
$\dd \logdet H = \tr(H^{-1} \dd H)$, together with the symmetry of $H$, one obtains
\begin{align}
\label{eq:gradA}
  \frac{\partial \Phi_{\mathrm{A}}}{\partial p}
  &= -2\, \tr\!\Big( H^{-2} J^\T \Gnoise^{-1} \frac{\partial J}{\partial p} \Big)
   = -2\, \Big\langle W_{\mathrm{A}},\; \frac{\partial J}{\partial p} \Big\rangle_F,
  & W_{\mathrm{A}} &= \Gnoise^{-1} J H^{-2}, \\
\label{eq:gradD}
  \frac{\partial \Phi_{\mathrm{D}}}{\partial p}
  &= -2\, \tr\!\Big( H^{-1} J^\T \Gnoise^{-1} \frac{\partial J}{\partial p} \Big)
   = -2\, \Big\langle W_{\mathrm{D}},\; \frac{\partial J}{\partial p} \Big\rangle_F,
  & W_{\mathrm{D}} &= \Gnoise^{-1} J H^{-1},
\end{align}
where $\langle A, B\rangle_F = \sum_{ij} A_{ij} B_{ij}$ is the Frobenius inner product. The
matrices $H^{-1}$ and $H^{-2}$ are computed by repeated triangular solves with the Cholesky
factor $L$ (functions \code{compute\_cost\_function\_gradient\_\{a,d\}\_opt\_optimized}).

Since only the rows of $J$ belonging to sensor $m$ depend on $\sensor{m}$, the Frobenius product
in \eqref{eq:gradA}--\eqref{eq:gradD} reduces to a sum over that sensor's rows:
\begin{equation}
  \frac{\partial \Phi}{\partial (\sensor{m})_p}
  \;=\;
  -2 \sum_{l=1}^{N} \sum_{k=1}^{n_{\mathrm{el}}}
     W_{(m,l),k}\; \frac{\partial J_{(m,l),k}}{\partial (\sensor{m})_p},
\end{equation}
which is exactly the loop structure of the gradient routines. For 3-axis MDEIT the same
expression is summed over the three axis blocks of the stacked Jacobian
$J = [J_x;\, J_y;\, J_z]$.

\section{The sensitivity Jacobian and its position derivatives}

\subsection{Jacobian with respect to conductivity (adjoint method)}

The datum for sensor $m$, stimulation $l$ (one fixed axis $d$; the axis index is suppressed) is
$B_m^{(l)} = (\Gamma\, u^{(l)})_m$ with $S(\sigma)\, u^{(l)} = b^{(l)}$. Differentiating with
respect to $\sigma_k$ and using
$\partial u^{(l)}/\partial \sigma_k = -S^{-1} (\partial S/\partial \sigma_k)\, u^{(l)}$ gives
\begin{equation}
  J_{(m,l),k}
  \;=\;
  \underbrace{\lambda_m^\T \frac{\partial S}{\partial \sigma_k}\, u^{(l)}}_{\text{adjoint part}}
  \;+\;
  \underbrace{\Big( \frac{\partial \Gamma}{\partial \sigma_k}\, u^{(l)} \Big)_{\!m}}_{\text{direct part}},
  \qquad
  S\, \lambda_m \;=\; -\,\Gamma^\T e_m ,
\end{equation}
where $e_m$ is the $m$-th unit vector; the adjoint solves are performed with preconditioned
conjugate gradients. By \eqref{eq:dSdsigma}, the adjoint part is
$|T_k|\, (\nabla \lambda_m)_k \cdot (\nabla u^{(l)})_k$, and since $\Gamma$ is linear in
$\sigma$, the direct part follows from \eqref{eq:Cmatrices} by replacing $\Sigma$ with
$e_k e_k^\T$. In the code these are the \code{dfdx} and \code{dfdp} terms of
\code{calc\_jacobian\_1axis\_direct\_fully\_vectorized\_local}.

\subsection{Position derivative of the Biot--Savart matrices}

The only quantities depending on $\sensor{m}$ are the rows of $R_x, R_y, R_z$. With
$\bm{d} = \sensor{m} - \bm{r}$, differentiating \eqref{eq:Rdef} under the integral sign gives
\begin{equation}
\label{eq:dRdp}
  \frac{\partial (R_j)_{mk}}{\partial (\sensor{m})_p}
  \;=\;
  \int_{T_k} \frac{\delta_{jp}\, |\bm{d}|^2 - 3\, d_j\, d_p}{|\bm{d}|^5}\, \dd V ,
  \qquad j,p \in \{x,y,z\},
\end{equation}
evaluated with the same quadrature rule (\code{dRmkj\_dp}, \code{dRmkj\_xyz}). This explicit,
smooth derivative is what replaces the electrode shape derivatives of the EIT setting.

\subsection{Position derivative of the Jacobian}

Differentiating $J_{(m,l),k}$ with respect to $p = (\sensor{m})_p$ produces two contributions
(functions \code{compute\_dJ\_ds}, \code{compute\_dJ\_xyz}, \code{compute\_dJxyz\_xyz}):
\begin{equation}
\label{eq:dJ}
  \frac{\partial J_{(m,l),k}}{\partial p}
  \;=\;
  \underbrace{|T_k| \,
    \Big(\nabla \frac{\partial \lambda_m}{\partial p}\Big)_{\!k} \cdot \big(\nabla u^{(l)}\big)_k}_{\code{dJ1}}
  \;+\;
  \underbrace{\frac{\partial}{\partial p}
    \Big( \frac{\partial \Gamma}{\partial \sigma_k}\, u^{(l)} \Big)_{\!m}}_{-\,\code{dJ2}} .
\end{equation}
\begin{itemize}
  \item The second term is obtained directly from \eqref{eq:Cmatrices} and \eqref{eq:dRdp}: it is
        the direct part of the Jacobian with $R_j$ replaced by $\partial R_j / \partial p$.
  \item The first term requires the derivative of the adjoint solutions. Since $S$ does not
        depend on the sensor positions, differentiating $S\lambda_m = -\Gamma^\T e_m$ gives one
        extra adjoint solve per sensor and coordinate,
        \begin{equation}
          S\, \frac{\partial \lambda_m}{\partial p}
          \;=\;
          -\, \frac{\partial \Gamma^\T}{\partial p}\, e_m ,
        \end{equation}
        where $\partial \Gamma / \partial p$ again follows from \eqref{eq:Gamma} and
        \eqref{eq:dRdp} (functions \code{compute\_dlambda\_ds}, \code{compute\_dlambda\_xyz},
        \code{compute\_dlambdaxyz\_xyz}).
\end{itemize}
Note that the forward solutions $u^{(l)}$ do not depend on the sensor positions at all --- another
consequence of the electrodes being fixed.

\section{Parametrization, constraints and optimization}

\subsection{Cylindrical coordinates and chain rule}

For a cylindrical tank it is convenient to parametrize each sensor by
$q_m = (r_m, \theta_m, z_m)$ with
$\bm{x}_m = (r_m \cos\theta_m,\, r_m \sin\theta_m,\, z_m)$. The gradient with respect to $q$
follows from the chain rule,
\begin{equation}
\label{eq:chainrule}
  \frac{\partial \Phi}{\partial q_i}
  \;=\;
  \sum_{p=1}^{3} \frac{\partial x_p}{\partial q_i}\, \frac{\partial \Phi}{\partial x_p},
  \qquad
  \frac{\partial \bm{x}}{\partial (r,\theta,z)}
  =
  \begin{pmatrix}
    \cos\theta & \sin\theta & 0 \\
    -r\sin\theta & r\cos\theta & 0 \\
    0 & 0 & 1
  \end{pmatrix},
\end{equation}
implemented in \code{dphidp\_to\_dphidq} with the transformation matrix from
\code{compute\_jacobian\_coordinate\_transformation\_cylindrical}.

\begin{remark}
The transformation matrix in \eqref{eq:chainrule} depends on the \emph{current} sensor
coordinates $(r_m, \theta_m)$ and must therefore be re-evaluated at every iterate of the
optimizer, not only at the initial configuration.
\end{remark}

Two constrained search domains are used with the interior-point algorithm of \code{fmincon}
(analytic gradients supplied, L-BFGS Hessian approximation):
\begin{itemize}
  \item \emph{cylinder shell}: $r_m = r_0$ fixed, $\theta_m$ free, $z_m \in [z_{\min}, z_{\max}]$;
  \item \emph{annular region}: $r_m \in [r_{\min}, r_{\max}]$, $\theta_m$ free,
        $z_m \in [z_{\min}, z_{\max}]$.
\end{itemize}
For unconstrained optimizers, bounded coordinates are mapped to $\R$ by the logistic
reparametrization $r = r_{\min} + (r_{\max}-r_{\min})\,\varsigma(\xi)$,
$\varsigma(\xi) = (1+e^{-\xi})^{-1}$, and analogously for $z$
(\code{map\_x\_to\_q\_cyl\_region} / \code{map\_q\_to\_x\_cyl\_region}).

\subsection{Sensor repulsion penalty}

Analogous to the electrode-gap penalty of Hyv\"onen et al., a smooth pairwise repulsion term can
be added to prevent sensor clustering:
\begin{equation}
  \Phi_{\mathrm{rep}}(\bm{s})
  \;=\;
  \frac{1}{2} \sum_{i \neq j} \frac{1}{|\bm{x}_i - \bm{x}_j|^2 + \epsilon},
  \qquad
  \frac{\partial \Phi_{\mathrm{rep}}}{\partial \bm{x}_i}
  \;=\;
  -2 \sum_{j \neq i} \frac{\bm{x}_i - \bm{x}_j}{\big(|\bm{x}_i - \bm{x}_j|^2 + \epsilon\big)^2},
\end{equation}
with a small $\epsilon > 0$ for smoothness (\code{repulsion\_cost},
\code{repulsion\_grad\_cartesian}). The total objective is
$\Phi = \Phi_{\mathrm{OED}} + \alpha\, \Phi_{\mathrm{rep}}$ with a manually chosen weight
$\alpha \geq 0$.

\section{Practical choices of the noise and prior covariances}
\label{sec:practical}

The script uses white (uncorrelated) covariances,
\begin{equation}
  \Gnoise = \beta\, I, \qquad \Gprior = \alpha_{\mathrm{pr}}\, I,
\end{equation}
scaled as follows:
\begin{itemize}
  \item The noise standard deviation is tied to the magnitude of the simulated data,
        $\sqrt{\beta} = \max_{j,m,l} |B_{j,m}^{(l)}| \,/\, 10$.
  \item The prior variance is set relative to the largest eigenvalue of $J_0^\T J_0$ (the
        Jacobian at the initial configuration),
        \begin{equation}
          \alpha_{\mathrm{pr}} \;=\; c\, \frac{\beta}{\lambda_{\max}(J_0^\T J_0)},
          \qquad c = 50 .
        \end{equation}
        The ratio $\alpha_{\mathrm{pr}}/\beta$ controls how many eigenmodes of the problem are
        \emph{data-dominated}: the number of eigenvalues $\lambda_i$ of $J_0^\T J_0$ with
        $\lambda_i\, \alpha_{\mathrm{pr}} / \beta > 1$ counts the modes for which the
        (linearized) data constrain the posterior more strongly than the prior. If this count is
        very small, the design criteria are dominated by the prior and become insensitive to the
        sensor positions; if it is very large, the linearization around the prior mean may become
        the limiting approximation.
\end{itemize}

\section{Summary of the algorithm}

\begin{enumerate}
  \item Build the FE model, fix electrodes, stimulation patterns, $\Gprior$, $\Gnoise$; draw an
        initial sensor configuration $\bm{s}^{(0)}$.
  \item At the current configuration: assemble $R_j(\bm{s})$, $\Gamma(\sigma_*; \bm{s})$; solve
        the forward problems $S u^{(l)} = b^{(l)}$ and adjoint problems
        $S \lambda_m = -\Gamma^\T e_m$; form the Jacobian $J$.
  \item Form $H = J^\T \Gnoise^{-1} J + \Gprior^{-1}$, its Cholesky factor, and evaluate
        $\Phi_{\mathrm{A}} = \tr(H^{-1})$ or $\Phi_{\mathrm{D}} = -\logdet H$.
  \item Compute $\partial R_j / \partial p$ from \eqref{eq:dRdp}, the extra adjoint solves for
        $\partial \lambda_m / \partial p$, and assemble $\partial J / \partial p$ from
        \eqref{eq:dJ}; contract with $W_{\mathrm{A}}$ or $W_{\mathrm{D}}$ as in
        \eqref{eq:gradA}--\eqref{eq:gradD}; transform the gradient to the optimization
        coordinates via \eqref{eq:chainrule}.
  \item Feed cost and gradient to the constrained optimizer (\code{fmincon}, interior-point,
        L-BFGS); iterate until convergence.
\end{enumerate}

\section{Case study: sensors confined to a circle at fixed height}
\label{sec:case-study}

This section documents a numerical case study, implemented in the sibling script
\code{example\_anomaly\_circle.m}, that specializes the theory above to the simplest
non-trivial circle-constrained design: $M$ magnetometers on a circle of \emph{fixed} radius
$r_m \equiv r_s$ and \emph{fixed} height $z_m \equiv z_s$ (half the tank height), so that only the
azimuths $\theta_m$ are free. Three-axis MDEIT is used throughout, $J = [J_x;J_y;J_z]$, and the
cylindrical chain rule \eqref{eq:chainrule} collapses to the one-dimensional form
\begin{equation}
\label{eq:chainrule1d}
  \frac{\partial \Phi}{\partial \theta_m}
  \;=\;
  -r_s \sin\theta_m\, \frac{\partial \Phi}{\partial x_m}
  \;+\;
  r_s \cos\theta_m\, \frac{\partial \Phi}{\partial y_m} ,
\end{equation}
evaluated at the \emph{current} $\theta_m$ at every iterate (cf.\ the remark in
\S\ref{sec:practical} above about the analogous $(r,\theta)$ dependence).

\subsection{Setup: ROI-informative prior and target anomaly}

A single conductive spherical anomaly (radius $R/3$, conductivity $5\times$ background) is placed
off-centre at half height, at $c_0 = (R/2,\,0,\,H/2)$, where $R,H$ are the tank radius and height.
The prior is informative about \emph{where} the anomaly may be but not its value, following the
``Case 1'' prior of Hyv\"onen et al.:
\begin{equation}
\label{eq:roiprior}
  (\Gprior)_{kk} \;=\;
  \begin{cases}
    \sigma_{\mathrm{roi}}^2, & \operatorname{centroid}(T_k) \in \mathcal{R}, \\[2pt]
    \sigma_{\mathrm{bg}}^2, & \text{otherwise},
  \end{cases}
  \qquad
  \mathcal{R} \;=\; \{\, k : |\operatorname{centroid}(T_k) - c_0| \le \rho \,\},
\end{equation}
with $\rho$ a multiple (\code{roi\_radius\_factor}) of the anomaly radius. Sensors are confined to
$r_m = r_s = 1.5R$, $z_m = z_s = H/2$ ($M = 8$), so \eqref{eq:chainrule1d} is the only
coordinate transformation needed.

\subsection{Noise calibration by target data-dominated rank}

Rather than the fixed $\sqrt\beta = \max|B|/10$ rule of \S\ref{sec:practical}, the case study
calibrates the noise level directly from the whitened-Jacobian spectrum at the even-spacing
configuration $\bm s_0$. Let $J_0 = J(\sigma_*;\bm s_0)$ and let $s_1 \ge s_2 \ge \dots$ be the
singular values of $J_0\, \Gprior^{1/2}$. Setting
\begin{equation}
\label{eq:noisecal}
  \sqrt\beta \;=\; s_{d}
\end{equation}
for a target rank $d$ (\code{d\_target}) makes exactly $d$ modes data-dominated
($s_i^2/\beta > 1$), placing the design problem in the regime where the sensor positions can
actually matter: too much noise and the criteria are prior-dominated and design-insensitive; too
little and every configuration already saturates them (cf.\ the discussion in
\S\ref{sec:practical}). A companion diagnostic, the mean ROI energy of the leading right
singular vectors $v_1,\dots,v_d$,
\begin{equation}
\label{eq:roienergy}
  \mathcal{E}_{\mathcal R} \;=\; \frac{1}{d} \sum_{i=1}^{d} \sum_{k \in \mathcal R} (v_i)_k^2 ,
\end{equation}
measures how much of the data-dominated subspace actually overlaps the ROI, as opposed to the
much larger and better-conditioned background.

\subsection{Numerical results}

With $\sigma_{\mathrm{bg}} = 0.005\,\sigma_0$, $\sigma_{\mathrm{roi}} = \sigma_0$
($\sigma_0$ the background conductivity), $\rho$ equal to one anomaly radius, and $d = 100$
($n_{\mathrm{el}} = 3425$, $n_{\mathrm{data}} = 384$), the even-spacing configuration already
reduces the ROI posterior variance by about $30\%$ relative to the prior
($\mathcal E_{\mathcal R} \approx 0.54$). Optimizing from this starting point (40 quasi-Newton
iterations of \code{fminunc}, analytic gradient, \code{SpecifyObjectiveGradient}) gives:
\begin{center}
\begin{tabular}{lccc}
\hline
Criterion & $\Phi(\text{even})$ & $\Phi(\text{optimized})$ & Gain \\
\hline
A-optimality, $\tr(H^{-1})$ & $2.879 \times 10^{-1}$ & $2.747 \times 10^{-1}$ & $4.6\%$ trace reduction \\
D-optimality, $-\logdet H$ & $-5.555 \times 10^{4}$ & $-5.557 \times 10^{4}$ & $17.6$ nats ($25.4$ bits) \\
\hline
\end{tabular}
\end{center}
Both criteria move the sensors toward the anomaly's azimuth, visibly more so for D-optimality. A
multistart check (3 random initial azimuth sets) converges to the same A-optimal value in every
run, and to within $0.03\%$ of the same D-optimal value, confirming that the even-spacing start is
already close to the global optimum of this low-dimensional ($M=8$) design space.

\subsection{Discussion: why the A-optimal gain stays modest}

\begin{remark}
Enlarging the data budget alone does not close the gap between the two criteria: repeating the
optimization for $d \in \{60, 100, 150\}$ leaves the converged A-optimal trace reduction in the
same $4$--$6\%$ range throughout, which rules out an insufficient-SNR explanation. Moving the
anomaly closer to the tank wall (from $R/2$ to $0.6R$, all else unchanged) leaves the (3-iterate)
A-optimal reduction almost unchanged ($2.7\% \to 4.4\%$) while nearly doubling the D-optimal
information gain ($13.8 \to 24.8$ nats). The bottleneck is therefore geometric, not statistical.
\end{remark}
A-optimality is an \emph{average} of marginal posterior variances,
$\Phi_{\mathrm{A}} = \sum_k (\Gpost)_{kk}$, and is comparatively insensitive to
\emph{correlations} between sensor measurements: two sensors whose Jacobian rows are nearly
collinear (typical for evenly spaced sensors seeing a common, roughly axisymmetric
electrode/stimulation pattern) each still contribute their own, largely redundant, reduction to
the trace. D-optimality, $\Phi_{\mathrm{D}} = -\logdet H = -\sum_i \log \lambda_i(H)$, behaves
oppositely: an eigenvalue of $H$ associated with a direction seen by two collinear sensors stays
pinned near its prior value, so \emph{removing} that redundancy by reallocating sensors is
disproportionately rewarded. This is consistent with the classical design-of-experiments folklore
that A-optimal designs tolerate point replication far better than D-optimal designs. Because the
fixed-radius, fixed-height circle used here is itself a highly symmetric search domain, the
even-spacing configuration is close to a ``replicated'' design in this sense: moving the sensors
mainly re-shuffles \emph{which} directions are well resolved (a large D-optimal effect) rather
than \emph{how much} variance is removed \emph{on average} (a small A-optimal effect). We have not
proved this as a theorem; it is offered as the most parsimonious explanation consistent with all
of the numerical evidence above. A natural follow-up, left to future work, is whether relaxing the
height constraint --- letting $z_m$ vary together with $\theta_m$, still at fixed radius $r_s$ ---
breaks enough of this residual symmetry to produce a more considerable A-optimal gain.

\subsection{Follow-up: freeing the height $z_m$}
\label{sec:case-study-thetaz}

The question closing the previous subsection is answered here using a sibling script,
\code{example\_anomaly\_circle\_thetaz.m}, identical to \code{example\_anomaly\_circle.m} in every
other setting (anomaly, prior, noise calibration, $d=100$, 3-axis MDEIT, \code{fminunc}
quasi-Newton with analytic gradient) except that each sensor now has \emph{two} free coordinates
$(\theta_m, z_m)$ at the same fixed radius $r_s = 1.5R$. To keep the unconstrained solver
\code{fminunc} usable with a single analytic gradient, $z_m$ is generated from an unconstrained
$\eta_m$ through a sigmoid,
\begin{equation}
\label{eq:sigmoidmap}
  z_m \;=\; z_{\min} + (z_{\max}-z_{\min})\,\mathrm{sigm}(\eta_m), \qquad
  \mathrm{sigm}(y) = \frac{1}{1+e^{-y}},
\end{equation}
with $z_{\min}=0$, $z_{\max}=H$, so that $\eta_m = 0$ reproduces the previous fixed height
$z_s = H/2$ exactly and $q_0 = (\theta_{\mathrm{even}}, \bm 0)$ is an identical-to-baseline start.
The chain rule \eqref{eq:chainrule1d} for $\theta_m$ is unchanged; height contributes the new term
\begin{equation}
\label{eq:etachainrule}
  \frac{\partial \Phi}{\partial \eta_m}
  \;=\;
  \frac{\partial \Phi}{\partial z_m}\,(z_{\max}-z_{\min})\,
  \mathrm{sigm}(\eta_m)\bigl(1-\mathrm{sigm}(\eta_m)\bigr),
\end{equation}
reusing the $p=3$ (Cartesian $z$) component of $\partial\Phi/\partial \bm p_m$ that the fixed-height
script already computed and discarded. Both halves of the resulting $2M$-length gradient were
verified against central finite differences (rel.\ err.\ $10^{-6}$--$10^{-8}$, the same standard
used throughout this study).

\paragraph{Results.} Optimizing from $q_0$ (40 quasi-Newton iterations, as in the theta-only case)
gives:
\begin{center}
\begin{tabular}{lccc}
\hline
Criterion & $\Phi(\text{even})$ & $\Phi(\text{optimized})$ & Gain \\
\hline
A-optimality, $\tr(H^{-1})$ & $2.879 \times 10^{-1}$ & $2.601 \times 10^{-1}$ & $9.7\%$ trace reduction \\
D-optimality, $-\logdet H$ & $-5.555 \times 10^{4}$ & $-5.558 \times 10^{4}$ & $29.0$ nats ($41.9$ bits) \\
\hline
\end{tabular}
\end{center}
A longer single-start run (150 iterations) confirms both optima are fully converged (A-optimal
gradient norm $7.6\times10^{-9}$; D-optimal run halts because it cannot decrease further along the
search direction, gradient norm $2.4\times10^{-5}$) at essentially the same values ($9.7\%$,
$29.0$ nats), so the 40-iteration figures above are not an artifact of premature stopping. A
multistart check (4 random initial $(\theta,\eta)$ sets, 60 iterations each) finds every run within
$0.3\%$ of the same A-optimal value and within $0.003\%$ of the same D-optimal value (best of the
four: $9.9\%$ trace reduction, $30.3$ nats), confirming that --- as in the theta-only case --- the
even-spacing start is close to the global optimum of this now $2M=16$-dimensional design space and
that no qualitatively better configuration is being missed.

\paragraph{Diagnosis.} Freeing $z_m$ roughly \emph{doubles} both gains relative to the
theta-only results of \S\ref{sec:case-study} ($4.6\%\to 9.7$--$9.9\%$ for A-optimality, $17.6\to
29$--$30$ nats for D-optimality), which is a real and reproducible effect, not noise: a non-trivial part of the
modest A-optimal gain in \S\ref{sec:case-study} \emph{was} a symmetry artifact of the
overly restrictive fixed-radius-fixed-height search domain, exactly as conjectured. However, the new
A-optimal reduction ($\approx 10\%$) still falls short of being ``considerable'' by the same
standard applied throughout this study ($\gtrsim 15$--$30\%$), so the restrictive-domain hypothesis
is only \emph{partially} confirmed. Inspecting the optimized $(\theta_m,z_m)$ pairs shows why:
sensors still cluster in azimuth toward the anomaly's side (as before), but in height they
\emph{spread out} across a substantial fraction of the shell ($\mathrm{std}(z_m) \approx 0.36$--$0.40\,$m
for A-optimality, $\approx 0.21\,$m for D-optimality, out of a $1.75\,$m range) rather than
converging tightly onto the anomaly's own height $z=H/2$. This is the same
decorrelation/de-clustering mechanism identified in \S\ref{sec:case-study} for azimuth, now acting
in the height direction as well: spreading in $z$ reduces redundancy between sensor rows (helping
both criteria, but disproportionately D-optimality) rather than concentrating averaged sensitivity
near the anomaly (which is what would move the A-optimal trace substantially). In short: the
fixed-height restriction explains \emph{part}, but evidently not \emph{all}, of the gap between the
two criteria's gains --- the remainder is consistent with the more fundamental explanation from
\S\ref{sec:case-study}, that A-optimality (an averaged criterion) is intrinsically more tolerant of
replicated/correlated sensor information than D-optimality, independent of how many geometric
degrees of freedom the design domain provides.

\section{Case study: human thorax, sensors on a fixed circle at half height}
\label{sec:case-study-thorax}

The two cylinder studies above used a rotationally symmetric conductor. This section repeats the
azimuth-only design (\S\ref{sec:case-study}) on a markedly more complicated geometry: an
adult-male thorax cross-section from the EIDORS \code{shape\_library}, extruded to a 3D slab of
height $H=0.5$, with the two lungs as interior low-conductivity regions ($\sigma_{\mathrm{bg}}=1$,
$\sigma_{\mathrm{lung}}=0.3$; this homogeneous-lung image is the OED linearization point / prior
mean). Twelve CEM electrodes form one ring at half height; $M=16$ magnetometers lie on a circle of
\emph{fixed} radius $R_0=1.5$ at $z_s=H/2$, only their azimuths $\theta_m$ free, so
\eqref{eq:chainrule1d} is again the only coordinate transformation. The implementation is
\code{human\_thorax\_theta.m}; it drives the shared, finite-difference-validated optimizer
\code{optimize\_sensor\_configuration.m} with the theta coordinate maps, and the analytic
$\partial\Phi/\partial\theta_m$ was re-verified against central finite differences (rel.\ err.\
$10^{-9}$--$10^{-6}$) at both the even start and a random perturbed configuration, for A- and
D-optimality.

\paragraph{Prior, noise, baselines.} The diagonal ``Case~1'' prior \eqref{eq:roiprior} is used,
with the suspected-anomaly centre $c_0=(0.1,0.1,H/2)$, $\sigma_{\mathrm{bg}}=0.005$ and
$\sigma_{\mathrm{roi}}=1$. On the coarse thorax mesh ($n_{\mathrm{el}}=5955$) a factor-one ROI
ball holds only $14$ elements, making the ROI-energy diagnostic \eqref{eq:roienergy}
pathologically small ($\mathcal E_{\mathcal R}\approx0.11$); enlarging the ball to
$\rho=2\times$ the anomaly radius restores a $129$-element ROI (comparable to the cylinder's
$\sim\!100$) with $\mathcal E_{\mathcal R}\approx0.33$. The noise is calibrated by
\eqref{eq:noisecal} with $d=100$ ($n_{\mathrm{data}}=576$), giving $\sqrt\beta\approx1.4\times
10^{-6}$; at this level the even configuration already reduces the ROI posterior variance by about
$26\%$ relative to the prior. Two baselines are compared against, both at $M=16$: (i) even
azimuthal spacing (also the optimizer start), and (ii) a ``boundary'' configuration with sensors
at a fixed standoff $\delta=0.05$ from the body along equiangular rays --- the thorax analogue of
``as close to the conductor as allowed''.

\paragraph{Results.} Optimizing from the even start (40 quasi-Newton iterations):
\begin{center}
\begin{tabular}{lccc}
\hline
Criterion & $\Phi(\text{even})$ & $\Phi(\text{optimized})$ & Gain vs.\ even \\
\hline
A-optimality, $\tr(H^{-1})$ & $9.5559\times10^{1}$ & $9.5415\times10^{1}$ & $0.15\%$ trace reduction \\
D-optimality, $-\logdet H$ & $-6.20443\times10^{4}$ & $-6.20482\times10^{4}$ & $3.9$ nats ($5.6$ bits) \\
\hline
\end{tabular}
\end{center}
The A-optimal even start reaches ``Local minimum found''; a multistart check (three random initial
azimuth sets) agrees to $\le\!0.003\%$ on the A-optimal value and to all printed digits on the
D-optimal value, so these are genuine converged optima, not premature stops. Crucially, the
``as close as allowed'' boundary configuration is \emph{far} better than any fixed-$R_0$ circle:
$\Phi_{\mathrm A}=71.08$ (a further $25\%$ below the azimuth optimum) and
$\Phi_{\mathrm D}=-6.312\times10^{4}$ (a further $1076$ nats), consistent with the $1/r^2$
Biot--Savart decay --- proximity, not azimuth, is the dominant lever, and it is entirely
unavailable at fixed $R_0$.

\paragraph{Discussion: weaker, not stronger, than the cylinder.} The a~priori expectation was that
the thorax, breaking rotational symmetry twice (non-circular boundary; off-centre ROI), would
yield a \emph{larger} A-optimal azimuth gain than the cylinder's $4.6\%$. The opposite occurs
($0.15\%$), and the reason is geometric and quantifiable. What the averaged A-optimal trace
responds to is the azimuthal \emph{contrast} in the ring's sensitivity to the ROI, i.e.\ the spread
of ROI-to-sensor distances around the circle. For the specified centre $c_0$ (radius
$|c_0|\approx0.14$ from the axis) against $R_0=1.5$, that distance varies only by $\pm9.4\%$ around
the ring; the cylinder's anomaly at $0.5R$ against a $1.5R$ ring varied by $\pm32.5\%$. The ROI
here is in effect \emph{more central} relative to the sensor circle than the cylinder anomaly was,
so the fixed-radius circle conserves the average sensitivity budget under azimuth moves even more
nearly than in \S\ref{sec:case-study}. The body's non-circular boundary does not rescue this: with
the sensors held far out at $R_0=1.5$, the anomaly signal $\delta B$ is governed by the
sensor-to-ROI distance, not by the (strongly azimuth-dependent) sensor-to-surface standoff, so the
boundary shape barely enters the ROI sensitivity. The optimized azimuths are nonetheless still
weakly \emph{structured} in the expected direction --- for both criteria the number of sensors
within $\pm45^\circ$ of the ROI azimuth ($45^\circ$) rises from $3$ to $5$, with sensors bunching
on the ROI/near-body (north-east) side and thinning on the far side (A-optimality: $9$ of the $16$
sensors end up in the $0$--$144^\circ$ arc; largest single-sensor move $\approx36^\circ$) --- by
the same redundancy-reduction mechanism as before, and D-optimality again shows the larger
relative effect. But on this near-central-ROI, far-ring design the resultant clustering is mild
($|\!\sum e^{i\theta_m}|/M\approx0.17$) and the absolute gains are small for both criteria. The
lesson carried to the radius-free study (Task~B) is that on the thorax the decisive degree of
freedom is radial standoff, not azimuth: the boundary configuration's large advantage is entirely
a proximity effect that only a free (or noise-traded) radius can access.

\section*{Symbol--code correspondence}

\begin{center}
\begin{tabular}{ll}
\hline
Symbol / quantity & Code \\
\hline
$S(\sigma)$, eq.~\eqref{eq:schur} & \code{M(img,sigma)} \\
$R_x,R_y,R_z$, eq.~\eqref{eq:Rdef} & \code{compute\_r\_matrices\_local} \\
$\Gamma_d$, eq.~\eqref{eq:Gamma} & \code{compute\_gamma\_matrices\_local} (\code{img.Gamma1..3}) \\
$J$ (adjoint + direct part) & \code{calc\_jacobian\_1axis\_direct\_fully\_vectorized\_local} \\
$\lambda_m$ & \code{compute\_lambda\_x} / PCG solves inside the Jacobian routine \\
$\partial R_j / \partial p$, eq.~\eqref{eq:dRdp} & \code{dRmkj\_dp}, \code{dRmkj\_xyz} \\
$\partial \lambda_m / \partial p$ & \code{compute\_dlambda\_ds}, \code{compute\_dlambda\_xyz} \\
$\partial J / \partial p$, eq.~\eqref{eq:dJ} & \code{compute\_dJ\_ds}, \code{compute\_dJ\_xyz}, \code{compute\_dJxyz\_xyz} \\
$\Phi_{\mathrm{A}}$, $\Phi_{\mathrm{D}}$ & \code{compute\_cost\_function\_\{a,d\}\_opt(\_3\_axis)} \\
$W_{\mathrm{A}}$, $W_{\mathrm{D}}$, eqs.~\eqref{eq:gradA}--\eqref{eq:gradD} &
  \code{compute\_cost\_function\_gradient\_\{a,d\}\_opt\_optimized(\_3\_axis)} \\
Chain rule, eq.~\eqref{eq:chainrule} & \code{dphidp\_to\_dphidq} \\
$\Phi_{\mathrm{rep}}$ & \code{repulsion\_cost}, \code{repulsion\_grad\_cartesian} \\
\hline
\end{tabular}
\end{center}

\section*{Reference}

N.~Hyv\"onen, A.~Sepp\"anen, S.~Staboulis,
\emph{Optimizing electrode positions in electrical impedance tomography},
SIAM Journal on Applied Mathematics 74(6), 1831--1851 (2014). arXiv:1404.7300.

\end{document}